\begin{abstract}
Robot demonstration data is expensive to collect; egocentric human data is cheap and abundant.
Mixing the two is increasingly common in practice, yet the design rules are empirical folklore: reported substitution ratios of 4--5 ego-hours per teleoperation-hour, occasional negative transfer, and no principled answer to \emph{when} and \emph{how much} egocentric data helps a robot policy.
We give a first-principles account. We model the two domains as Gaussian covariates with power-law feature spectra, the embodiment gap as a task-vector mismatch decomposed in the shared spectral basis, and the policy fit as mixture-weighted ridge regression evaluated on robot-domain risk.
Our main theorem characterizes the mixture risk in closed per-coordinate form (Monte-Carlo verified to 1--2\% relative error) and reveals a three-term structure of the marginal benefit of human data: an alignment term, a variance-reduction term, and a spectral-coverage term.
From it we derive an exact phase-transition criterion: mixing is beneficial at all robot sample sizes iff a functional $H$ of the mismatch spectrum and the spectra is non-positive; otherwise a critical robot sample size $\nr^{\mathrm{crit}} = K/H$ appears beyond which human data should be discarded --- recovering the classical source-discard transition as regularization vanishes.
An exchange-rate corollary defines $\kex$, the number of robot demonstrations equivalent to one egocentric demonstration, shows it is increasing in the inter-domain spectral gap and capped at $1$, and lands in the same range as empirically reported substitution ratios.
We validate the structural predictions in a controlled dual-domain Push-T testbed with diffusion policies: the first batch confirms the predicted non-monotone benefit with an interior optimum (up to $+32\%$ closed-loop coverage) and harm boundary, shows the benefit grows with training budget and that alignment strength is itself non-monotone, and surfaces an open-loop/closed-loop divergence as a target for the next theory iteration.
\end{abstract}
